\documentclass[11pt,a4paper]{article}
\usepackage[T1]{fontenc}
\usepackage[utf8]{inputenc}
\usepackage{newtxtext}
\usepackage[scaled=.92]{inconsolata}
\usepackage[margin=25mm,headheight=15pt]{geometry}
\usepackage{microtype,amsmath,amsthm,mathtools,booktabs,array,longtable}
\usepackage{newtxmath}
\usepackage{xcolor,listings,enumitem,fancyhdr}
\usepackage[most]{tcolorbox}
\usepackage{hyperref}
\hypersetup{colorlinks=true,linkcolor=black,citecolor=black,urlcolor=blue!45!black,
 pdftitle={Decomposing Algebraic Roots into Low-Degree Sums and Products},
 pdfauthor={Technical study prepared for Vladimir Reshetnikov}}
\definecolor{ink}{RGB}{42,53,44}
\definecolor{olive}{RGB}{84,96,64}
\definecolor{paper}{RGB}{247,247,240}
\definecolor{rulecolor}{RGB}{192,197,174}
\lstdefinelanguage{Wolfram}{
 morekeywords={Module,With,Block,If,Do,While,Table,Function,True,False,Automatic,
 All,Root,RootReduce,MinimalPolynomial,Resultant,ToNumberField,
 AlgebraicNumber,AlgebraicNumberPolynomial,PolynomialRemainder,PolynomialGCD,
 CoefficientList,NullSpace,LinearSolve,Failure,BeginPackage,EndPackage,
 Begin,End,SetAttributes,HoldAll,Catch,Throw,Options,OptionsPattern,OptionValue},
 sensitive=true,morecomment=[s]{(*}{*)},morestring=[b]"}
\lstset{language=Wolfram,basicstyle=\ttfamily\footnotesize,
 keywordstyle=\color{ink}\bfseries,commentstyle=\color{olive},stringstyle=\color{ink},
 backgroundcolor=\color{paper},frame=single,rulecolor=\color{rulecolor},
 breaklines=true,breakatwhitespace=false,columns=fullflexible,keepspaces=true,
 showstringspaces=false,tabsize=2,aboveskip=9pt,belowskip=9pt,
 xleftmargin=3pt,xrightmargin=3pt,framesep=5pt}
\newtcolorbox{keybox}{colback=paper,colframe=rulecolor,boxrule=.6pt,
 arc=0pt,left=10pt,right=10pt,top=7pt,bottom=7pt,breakable}
\newtheorem{theorem}{Theorem}[section]
\newtheorem{proposition}[theorem]{Proposition}
\newtheorem{lemma}[theorem]{Lemma}
\newtheorem{corollary}[theorem]{Corollary}
\theoremstyle{definition}\newtheorem{definition}[theorem]{Definition}
\theoremstyle{remark}\newtheorem{remark}[theorem]{Remark}
\newcommand{\Q}{\mathbb Q}
\newcommand{\C}{\mathbb C}
\newcommand{\Z}{\mathbb Z}
\newcommand{\Gal}{\operatorname{Gal}}
\newcommand{\Tr}{\operatorname{Tr}}
\newcommand{\Nm}{\operatorname{N}}
\newcommand{\Span}{\operatorname{span}_{\Q}}
\newcommand{\dega}{\operatorname{deg}_{\Q}}
\newcommand{\Dplus}{D_{+}}
\newcommand{\Dtimes}{D_{\times}}
\newcommand{\Dtwo}{D_{\times,2}}
\newcommand{\code}[1]{\texttt{\detokenize{#1}}}
\setlength{\parindent}{0pt}
\setlength{\parskip}{5pt plus 1pt minus 1pt}
\setlist{itemsep=2pt,topsep=4pt}
\allowdisplaybreaks[2]
\pagestyle{fancy}
\fancyhf{}\fancyhead[L]{\small\color{olive}Low-degree algebraic decompositions}
\fancyhead[R]{\small\thepage}
\renewcommand{\headrulewidth}{.3pt}
\begin{document}
\begin{titlepage}
\vspace*{10mm}
{\Large\color{olive}Exact algebraic computation}\par\vspace{8mm}
{\huge\bfseries Decomposing Algebraic Roots\par
into Low-Degree Sums\par and Products\par}
\vspace{7mm}
{\Large Mathematica algorithms, optimality certificates,\par
and the role of the splitting field}\par
\vspace{10mm}
{\large Technical study prepared for Vladimir Reshetnikov}\par
6 September 2026
\vspace{11mm}
\begin{abstract}
An algebraic number of degree nine can be a sum or a product of two cubic
numbers, but finding such an identity and proving that it minimizes the
largest component degree are different tasks. We give exact solutions and
optimality proofs for both examples in Mathematica StackExchange question
105933. We then develop a terminating mathematical algorithm for the
unrestricted additive problem, using fixed fields in a Galois closure, and
a terminating algorithm for the unrestricted \emph{two-factor} product
problem, using relative norms and rational linear algebra.

A concrete degree-eight example shows that optimal multiplicative factors
can lie outside even the input's splitting field. It also disproves a
published multiplicative descent assertion whose proof divides by a
possibly zero coefficient. The corrected two-factor criterion explicitly
retains a bounded radical exponent. For arbitrarily many factors we supply
exact, certificate-guided bounded search, not a claimed general terminating
minimizer. The accompanying Wolfram Language package implements these
procedures; independent exact Python checks verify the polynomial identities,
minimal degrees, and finite-group calculations.
\end{abstract}
\vfill
\begin{keybox}
\textbf{Scope and validation.} Both requested degree-nine examples are
settled globally: the optimal maximum degree is \textbf{three}, even with
arbitrarily many components allowed. General sums and general two-factor
products have complete finite algorithms here. The arbitrary-length product
problem is covered by certified search and lower bounds. Native Wolfram
execution could not be completed because the available evaluator returned
HTTP 404; the supplied native test suites are therefore unexecuted.
\end{keybox}
\end{titlepage}
{\small\setlength{\parskip}{0pt}\tableofcontents}
\newpage

\section{The problem, made precise}

For an exact algebraic number $a$, write
\[
 \dega(a)=[\Q(a):\Q].
\]
The degree is always the \emph{absolute degree over the rationals}, not the
degree of an arbitrary polynomial that happens to vanish at $a$, and not a
relative extension degree over a field already containing other components.
Define
\begin{align}
 \Dplus(a)&=\min\left\{\max_i\dega(b_i):
             a=b_1+\cdots+b_r,\quad r\ge1\right\},\label{eq:dplus}\\
 \Dtimes(a)&=\min\left\{\max_i\dega(b_i):
             a=b_1\cdots b_r,\quad r\ge1,\ b_i\ne0\right\},\label{eq:dtimes}
\end{align}
where the $b_i$ may be anywhere in $\overline\Q$. The multiplicative
definition initially assumes $a\ne0$; zero has the trivial degree-one
representation. Rationals have degree one. Permitting a single component
makes the feasible upper bound $\dega(a)$ explicit. Requiring at least two
components instead changes nothing, because zero or one can be appended.

The number of components is not the objective. Nor are coefficient height,
radical nesting, expression length, or numerical conditioning. Those are
useful secondary objectives, but optimizing them must not silently replace
\eqref{eq:dplus} or \eqref{eq:dtimes}.

We will also need the genuinely different quantity
\[
 \Dtwo(a)=\min_{a=bc}\max\{\dega(b),\dega(c)\}.
\]
Always $\Dtimes(a)\le\Dtwo(a)$, and strict inequality occurs. A complete
algorithm for two factors is consequently not, by itself, a complete
algorithm for the question as stated.

Three ambient domains must also be distinguished:
\[
 K=\Q(a)\ \subseteq\ L=\text{the splitting field of the minimal polynomial of }a
 \ \subseteq\ \overline\Q.
\]
Restricting components to $K$ can change the answer. For sums it is enough to
work in $L$, by a trace argument. For products even restriction to $L$ can
change the answer. These distinctions drive the implementation.

\section{The two examples: exact answers and short certificates}
\label{sec:examples}

Let $u$ and $v$ be the unique real roots of
\[
 p(x)=x^3+x+1,\qquad q(x)=x^3-x+1.
\]
Both cubics are irreducible over $\Q$ by the rational root test. Their
discriminants are $-31$ and $-23$, respectively, so each has one real root.
Put $a=uv$ and $s=u+v$. Their minimal polynomials are
\begin{align}
 P_\times(x)&=x^9+2x^7-3x^6+x^5-x^4+3x^3-x-1,\label{eq:productpoly}\\
 P_+(x)&=x^9+6x^6+3x^5-15x^3+24x^2-4x+8.\label{eq:sumpoly}
\end{align}
The exact verification script supplied with this article checks
irreducibility and counts the real roots of both polynomials; each has
exactly one. Thus the first real \code{Root} in each example is the intended
number.

\subsection{A compact Mathematica answer}

\begin{lstlisting}
p = x^3 + x + 1;
q = x^3 - x + 1;
u = Root[1 + # + #^3 &, 1];
v = Root[1 - # + #^3 &, 1];

ap = Root[-1 - # + 3 #^3 - #^4 + #^5 - 3 #^6 +
          2 #^7 + #^9 &, 1];
as = Root[8 - 4 # + 24 #^2 - 15 #^3 + 3 #^5 +
          6 #^6 + #^9 &, 1];

RootReduce /@ {ap - u v, as - u - v}
(* Expected: {0, 0} *)

Exponent[MinimalPolynomial[#, x], x] & /@ {ap, as, u, v}
(* Expected: {9, 9, 3, 3} *)
\end{lstlisting}
\code{RootReduce} performs exact algebraic normalization; it does not seek
this decomposition objective \cite{rootreduce}. Keeping the components in
a list, or displaying an inactive sum/product, avoids immediately collapsing
an answer back into a single algebraic number.

The expected display values, not used as proof, are
\[
 u\approx-0.68232780382801932737,\quad
 v\approx-1.32471795724474602596,
\]
\[
 a\approx0.90389189445834756110,\qquad
 s\approx-2.00704576107276535333.
\]

\subsection{Recovering the cubics directly from the degree-nine inputs}

The following rational functions are particularly short exact certificates:
\begin{equation}
 u=\frac{a^3-a^2-1}{a^4+a^2-a+1},\qquad v=\frac{a}{u},
 \label{eq:recoverproduct}
\end{equation}
and
\begin{equation}
 u=\frac{3s^5-5s^3-3s^2+2s-4}{6s^4-6s+4},\qquad v=s-u.
 \label{eq:recoversum}
\end{equation}
Their denominators are relatively prime to the respective input minimal
polynomials, so they do not vanish at any conjugate of the input. Substituting
these expressions into $p(u)$ and $q(v)$ gives zero modulo
$P_\times(a)$ or $P_+(s)$. Since all values in these formulas are real, the
unique real roots of the cubics are selected automatically.

For example, the product calculation can be reproduced entirely with
polynomial remainders:
\begin{lstlisting}
P = x^9 + 2 x^7 - 3 x^6 + x^5 - x^4 + 3 x^3 - x - 1;
num = x^3 - x^2 - 1;
den = x^4 + x^2 - x + 1;
{
 PolynomialGCD[P, den],
 PolynomialRemainder[num^3 + num den^2 + den^3, P, x],
 PolynomialGCD[P, num],
 PolynomialRemainder[(x den)^3 - (x den) num^2 + num^3, P, x]
}
(* Expected: {1, 0, 1, 0} *)
\end{lstlisting}
The first remainder certifies $p(\mathrm{num}/\mathrm{den})=0$; the second
certifies $q(x\,\mathrm{den}/\mathrm{num})=0$. No numerical root matching is
involved. Rational recovery formulas can be obtained by a polynomial GCD or
extended Euclidean computation over $\Q(a)$ after eliminating one variable.
They are compact certificates, not a general discovery algorithm.

An alternative product certificate involves no denominators depending on
$a$:
\begin{align*}
 u&=-\frac{a^7+2a^5-2a^4+1}{2},\\
 v&=-\frac{a^7+2a^5-2a^4+2a^3+2a-1}{2}.
\end{align*}
It also demonstrates explicitly that both cubic factors lie in $\Q(a)$ for
this particular example. That fortunate property is not a general theorem.

\subsection{Why degree three is optimal, not merely successful}

Suppose a number is a sum or product of finitely many numbers of degree at
most two. Each component lies in a quadratic or rational field. Their
compositum is a multiquadratic extension and has degree a power of two.
Every element of that compositum has degree dividing that power of two.
A degree-nine number cannot occur. Consequently both displayed
representations attain the global minimum:
\[
 \boxed{\Dtimes(a)=3,\qquad \Dplus(s)=3.}
\]
This rules out \emph{any number} of quadratic components, even in auxiliary
fields not suggested by the original expression. Section~\ref{sec:bounds}
generalizes this certificate.

\section{Discovering candidates with resultants and exact residuals}
\label{sec:discovery}

\subsection{Forward construction is easy}

If $p(u)=q(v)=0$, then an annihilating polynomial for $u+v$ is
\[
 R_+(x)=\operatorname{Res}_y\bigl(p(y),q(x-y)\bigr).
\]
If $m=\deg q$, an annihilating polynomial for $uv$ is
\[
 R_\times(x)=\operatorname{Res}_y\bigl(p(y),y^m q(x/y)\bigr).
\]
The second argument in the latter expression is polynomial after clearing
the displayed denominator. The companion package implements these as
\code{SumPolynomial} and \code{ProductPolynomial}, using the documented
\code{Resultant} operation \cite{resultant}.

For the examples:
\begin{lstlisting}
Resultant[y^3 + y + 1, (x-y)^3 - (x-y) + 1, y]
(* Pplus[x] *)
Resultant[y^3 + y + 1, x^3 - x y^2 + y^3, y]
(* Ptimes[x] *)
\end{lstlisting}
A characteristic polynomial of a Kronecker sum or Kronecker product of
companion matrices gives the same forward information. The existing
StackExchange answers use these sorts of constructions and coefficient
matching \cite{question}. They are useful discovery devices, but do not by
themselves prove minimality.

\subsection{Why blindly inverting a resultant is insufficient}

For unknown monic polynomials $p$ and $q$, one can equate resultant
coefficients to those of a target polynomial and solve the resulting system.
This is useful at very small degrees, especially after a plausible shape has
been identified. Several qualifications are essential.

First, the input minimal polynomial only has to \emph{divide} the
resultant. Equality is not generally necessary: repeated conjugate products
or sums can make the resultant a power of a smaller polynomial. In
Section~\ref{sec:external}, two quartics have a degree-eight product and the
resultant is the square of its minimal polynomial.

Second, unknown coefficients must be rational. A complex or algebraic
solution of the coefficient equations does not give the intended absolute
degree bound. Third, integer monic polynomials only cover algebraic
integers. A complete finite-height dictionary must include nonmonic
primitive integer minimal polynomials. Finally, convenient normalizations
such as constant term one or simultaneously depressed cubics need their own
justification; rational rescaling cannot always achieve them.

After any candidate polynomials have been found, enumerate their exact root
branches and verify the actual input value. A resultant identifies possible
conjugate values, not a particular branch. Mathematica places real polynomial
roots first in increasing order, but complex-root ordering has additional
method-dependent details \cite{root}. Do not reuse a numerically guessed
root index as a universal branch rule.

\subsection{The residual trick: search for one polynomial, not two}

When a candidate polynomial $p$ is available, enumerate its roots $b$ and
compute
\[
 c=a-b\quad\text{or}\quad c=a/b.
\]
Then calculate the exact minimal polynomial of $c$. The second unknown
polynomial does not need to be searched independently. This reduces a
candidate-pair problem to a candidate-singleton problem.

After loading the package, the examples are found by
\begin{lstlisting}
Get["code/RootDecomposition.wl"];
PairFromPolynomial[ap, x^3+x+1, x, "Product", 3]
PairFromPolynomial[as, x^3+x+1, x, "Sum", 3]
\end{lstlisting}
A successful association contains the component list, each absolute degree,
the maximum degree, exact equality status, and an optimality flag. A failed
trial means only that \emph{this candidate polynomial} did not produce a
pair under the degree cap.

Without advance knowledge of either cubic, the following bounded calls
search the primitive integer polynomials of degree at most three and
coefficient height at most one:
\begin{lstlisting}
SearchRootDecomposition[ap, "Product", 3, 1, 2]
SearchRootDecomposition[as, "Sum",     3, 1, 2]
\end{lstlisting}
Both cubics belong to that dictionary. The exact arithmetic checks and the
prime-divisor bound provide the certificates; dictionary membership is only
the discovery mechanism. The search implementation and its precise
completeness statement are discussed in Section~\ref{sec:many}.

\section{Lower bounds that survive arbitrary auxiliary fields}
\label{sec:bounds}

\begin{theorem}[Normal-closure obstruction]\label{thm:bound}
Let $a$ have degree $n$, let $L$ be its splitting field, and put
$G=\Gal(L/\Q)$. If $a$ is a finite sum or product of algebraic numbers of
degree at most $d$, then
\begin{equation}
 \text{every prime divisor of }|G|\text{ is at most }d,
 \qquad
 \exp(G)\mid\operatorname{lcm}(1,2,\ldots,d).
 \label{eq:galoisbound}
\end{equation}
In particular, the largest prime divisor of $n$ is a lower bound for both
$\Dplus(a)$ and $\Dtimes(a)$.
\end{theorem}
\begin{proof}
For every component $b_i$, let $M_i$ be the normal closure of $\Q(b_i)$.
If $m_i=\dega(b_i)$, the action on its $m_i$ conjugates embeds
$\Gal(M_i/\Q)$ in $S_{m_i}$. Let $M$ be the compositum of these normal
closures. Restriction embeds $\Gal(M/\Q)$ into their direct product. Since
$M/\Q$ is normal and contains $a$, it contains every conjugate of $a$, hence
$L$. Restriction onto $G$ is surjective.

All prime divisors of the order of a subgroup of the direct product are at
most $d$. Its exponent divides $\operatorname{lcm}(1,\ldots,d)$, because
that is a common multiple of the orders of all permutations in each
$S_{m_i}$. These properties pass to a quotient. Finally $n$ divides
$[L:\Q]=|G|$.
\end{proof}

\begin{corollary}
A number of prime degree $p$ cannot be expressed as a finite sum or product
of numbers all having degrees less than $p$.
\end{corollary}

The inexpensive implementation is
\begin{lstlisting}
PrimeDegreeBound[a]
\end{lstlisting}
which needs only \code{MinimalPolynomial} and integer factorization. Once
Galois data are available, the package strengthens it using the least $d$
for which the exponent divisibility in \eqref{eq:galoisbound} holds. For
example, exponent four rules out degree caps two \emph{and three}.

A separate, length-dependent bound is
\begin{equation}
 \dega(a)\le\prod_{i=1}^{r}\dega(b_i)\le d^r.
 \label{eq:lengthbound}
\end{equation}
It follows from the usual degree inequality for a compositum. For two
components it gives $d\ge\lceil\sqrt n\rceil$.

\begin{remark}
Do not strengthen the compositum inequality to an unjustified divisibility
assertion. For nonnormal fields, the compositum degree need not divide the
product of their degrees: two distinct conjugate cubic fields in an
$S_3$ splitting field can have compositum degree six, although their degree
product is nine. The prime-divisor proof above deliberately uses normal
closures, whose groups embed in symmetric groups.
\end{remark}

These bounds are necessary, not generally sufficient. A certificate saying
\code{GlobalOptimal -> False} in the lightweight checker means that global
optimality was \emph{not established}, not that the representation is known
to be suboptimal.

\section{A complete solution for arbitrary finite sums}
\label{sec:sums}

\subsection{Why all summands can be moved into a finite field}

The following is the decisive additive descent principle. An additive
normal-closure reduction also appears in Altamirano's study
\cite{altamirano}; we give the argument explicitly, including its
componentwise degree bound.

\begin{lemma}[Trace projection]\label{lem:trace}
Let $L/\Q$ be finite Galois, let $a\in L$, and suppose
$a=\sum_i b_i$ with arbitrary algebraic $b_i$. Then there are
$c_i\in L\cap\Q(b_i)$ with
\[
 a=\sum_i c_i,\qquad \dega(c_i)\le\dega(b_i).
\]
\end{lemma}
\begin{proof}
Choose a finite Galois extension $M/\Q$ containing $L$ and every $b_i$,
and put $N=\Gal(M/L)$. Define
\[
 \pi_L(x)=\frac1{|N|}\sum_{\sigma\in N}\sigma(x).
\]
This is a $\Q$-linear projection onto $L$, so
$a=\pi_L(a)=\sum_i\pi_L(b_i)$. It remains to control degree, rather than
merely field membership.

Because $L/\Q$ is Galois, $N$ is normal in $\Gal(M/\Q)$. If $h$ fixes
$b_i$, conjugation by $h$ permutes $N$, and therefore
$h\pi_L(b_i)=\pi_L(hb_i)=\pi_L(b_i)$. Thus $\pi_L(b_i)$ is fixed by
$\Gal(M/\Q(b_i))$ as well as by $N$. It lies in
$\Q(b_i)\cap L$. Set $c_i=\pi_L(b_i)$.
\end{proof}

The normality assumption is essential to the equivariance step. Averaging
onto the generally nonnormal field $\Q(a)$ need not preserve the component
degree. The correct finite ambient field is the splitting field $L$.

\subsection{Replace unknown algebraic numbers by fixed spaces}

Let $G=\Gal(L/\Q)$. For $H\le G$, the fixed field $L^H$ has absolute
degree $[G:H]$. Define the rational vector subspace
\begin{equation}
 W_d=\sum_{\substack{H\le G\\{}[G:H]\le d}} L^H\ \subseteq L.
 \label{eq:Wd}
\end{equation}
The sum here is a \emph{linear sum of subspaces}. It is not the compositum
of those fields. Passing to the compositum would allow products and would
solve a different problem.

\begin{theorem}[Exact additive criterion]\label{thm:sum}
For an algebraic number $a$ with splitting field $L$,
\[
 \Dplus(a)\le d\quad\Longleftrightarrow\quad a\in W_d.
\]
Consequently $\Dplus(a)$ is computable by finite subgroup enumeration and
rational linear algebra.
\end{theorem}
\begin{proof}
If a representation with maximum degree at most $d$ exists, apply
Lemma~\ref{lem:trace}. Each projected summand $c_i$ lies in a field
$\Q(c_i)\subseteq L$ of degree at most $d$, so it lies in one of the
summands in \eqref{eq:Wd}. Conversely, a vector-space representation in
$W_d$ is a finite sum of elements of the listed fields, each of degree at
most $d$.

There are finitely many subgroups of the finite group $G$. Test $d$ in
increasing order up to $n=\dega(a)$. The test must succeed by $d=n$,
because $a\in\Q(a)\subseteq L$.
\end{proof}

A rational scalar multiplying a basis vector can be absorbed into that
summand without increasing degree. Selecting an independent subset of the
combined field bases shows that no more than $[L:\Q]$ individual basis
contributions are ever needed. In practice, grouping all contributions from
the same fixed field often gives a much shorter answer.

\subsection{The exact matrix computation}

Choose a primitive element $\theta$ for $L$ and use the power basis
$1,\theta,\ldots,\theta^{m-1}$, where $m=[L:\Q]$. Represent field
elements as rational column vectors. For each $g\in G$, let $A_g$ be its
$m\times m$ rational matrix. A basis matrix for $L^H$ is obtained from
\[
 \bigcap_{h\in H}\ker(A_h-I).
\]
Stacking the equations and calling \code{NullSpace} computes this basis.
For a fixed degree cap, concatenate the eligible field bases into a matrix
$B_d$. Then
\[
 a\in W_d\quad\Longleftrightarrow\quad B_d c=\operatorname{vec}(a)
\]
has a rational solution. \code{LinearSolve} supplies one when it exists
\cite{linearsolve}. Its coefficients are grouped by field and converted back
into exact algebraic numbers.

The implementation stores fixed-field bases as \emph{rows}; it transposes
them when forming $B_d$. Automorphism and multiplication matrices act on
\emph{columns}. This convention matters: accidentally transposing one
operator can yield a plausible-looking but incorrect subfield calculation.

\begin{lstlisting}
g = BuildGaloisData[as];
r = CompleteSumDecomposition[as, g];
r["MaximumDegree"]       (* expected 3 *)
r["GlobalOptimal"]       (* expected True *)
\end{lstlisting}
The example does not require this expensive route because the cubic
candidate and the prime bound already settle it. The fixed-space method is
the general completeness mechanism.

\subsection{Why the procedure is finite but can be expensive}

A degree-$n$ polynomial may have splitting field degree as large as $n!$.
All computations above take place in that field, not necessarily in the
much smaller field $\Q(a)$. Exact coefficients can also grow considerably.
The supplied backend is an explicit small-field reference implementation,
not a claim of a polynomial-time simplifier.

The defaults limit the constructed field to degree 128 and subgroup
enumeration to 10,000 entries. Reaching either guard produces a
\code{Failure}, not an impossibility certificate. The field-degree guard
is checked \emph{after} primitive-element construction; it does not bound
the time spent constructing that element. An outer \code{TimeConstrained}
can impose an operational time limit, again with an inconclusive result.

\section{The product of two factors: a complete norm--intersection algorithm}
\label{sec:two}

Unlike additive projection, a field norm preserves a product but raises the
input to a power. That power must be retained. It is precisely what permits
optimal factors outside the splitting field.

\subsection{An intersection lemma for Galois base fields}

\begin{lemma}\label{lem:intersection}
Let $L/\Q$ be finite Galois and $b$ algebraic. Put
$E=\Q(b)$, $E_0=E\cap L$, and $t=[L(b):L]$. Then
\[
 [E:\Q]=t[E_0:\Q],\qquad
 \Nm_{L(b)/L}(b)\in E_0.
\]
The minimal polynomial of $b$ over $L$ is the same polynomial as its minimal
polynomial over $E_0$.
\end{lemma}
\begin{proof}
The compositum $LE/E$ is Galois. Restriction identifies its Galois group
with a subgroup of $\Gal(L/\Q)$ whose fixed field in $L$ is $L\cap E$.
Thus $[LE:E]=[L:E_0]$. The tower formula gives
$[LE:L]=[E:E_0]$. The minimal polynomial over $E_0$ has this latter
degree, remains a polynomial over $L$, and vanishes at $b$, so it is already
minimal over $L$. Its constant coefficient, up to sign, is the displayed
norm, and belongs to $E_0$.
\end{proof}

This standard norm/intersection fact is also stated in Samuels's
Lemma 2.2 \cite{samuels}. Its relevance here is the exact absolute-degree
factorization $t[E_0:\Q]$, not a height estimate.

\subsection{A finite necessary-and-sufficient criterion}

\begin{theorem}[Two-factor criterion]\label{thm:two}
Let $0\ne a\in L$, with $L/\Q$ finite Galois, and let $d\ge1$. There
exist $b,c\in\overline\Q$ such that
\[
 a=bc,\qquad \dega(b),\dega(c)\le d
\]
if and only if there are subfields $E_0,F_0\subseteq L$ and an integer
$t\ge1$ for which
\begin{equation}
 t[E_0:\Q]\le d,\qquad t[F_0:\Q]\le d,\qquad
 E_0\cap a^t F_0\ne\{0\}.
 \label{eq:paircriterion}
\end{equation}
If $n=\dega(a)$, it suffices to consider
\begin{equation}
 1\le t\le\min\left(d,\left\lfloor\frac{d^2}{n}\right\rfloor\right).
 \label{eq:tbound}
\end{equation}
\end{theorem}
\begin{proof}
For necessity, suppose $a=bc$. Then $L(b)=L(c)=M$, because
$c=a/b$ and $b=a/c$. Put $t=[M:L]$,
$E_0=L\cap\Q(b)$, and $F_0=L\cap\Q(c)$. By
Lemma~\ref{lem:intersection},
\[
 \dega(b)=t[E_0:\Q],\qquad\dega(c)=t[F_0:\Q].
\]
Let $u=\Nm_{M/L}(b)$ and $v=\Nm_{M/L}(c)$. These nonzero values belong
to $E_0$ and $F_0$, respectively, and
\[
 uv=\Nm_{M/L}(a)=a^t.
\]
Therefore $u=a^t/v\in E_0\cap a^tF_0$, establishing
\eqref{eq:paircriterion}.

For sufficiency, choose any nonzero $u\in E_0\cap a^tF_0$. Choose a
$t$-th root $b$ of $u$, and define $c=a/b$. Then
\[
 b^t=u\in E_0,\qquad c^t=a^t/u\in F_0.
\]
The polynomials $X^t-u$ over $E_0$ and $X^t-a^t/u$ over $F_0$ give
\[
 \dega(b)\le t[E_0:\Q]\le d,\qquad
 \dega(c)\le t[F_0:\Q]\le d.
\]
There is no root-of-unity ambiguity in the product: the second factor was
defined by division, rather than by an independent root choice.

Finally, writing $e=[E_0:\Q]$ and $f=[F_0:\Q]$, the condition implies
$a^t\in E_0F_0$. Consequently
\[
 n\le t\dega(a^t)\le t[E_0F_0:\Q]\le tef\le d^2/t.
\]
Thus $t\le d^2/n$, and $t\le d$ follows already from $te\le d$.
\end{proof}

The same proof works with unequal degree caps $d_1,d_2$, replacing the first
two inequalities by $te\le d_1$, $tf\le d_2$, and the upper bound by
$t\le\min(d_1,d_2,\lfloor d_1d_2/n\rfloor)$. The companion package uses
the equal-cap form because that is the requested objective.

\subsection{Everything before the final roots is rational linear algebra}

Let $B_E,B_F$ be column basis matrices of $E_0,F_0$ in a common power
basis of $L$. Let $M_{a^t}$ be the rational matrix of multiplication by
$a^t$. A nonzero element in the intersection is obtained from a nonzero
solution of
\begin{equation}
 \begin{bmatrix}B_E&-M_{a^t}B_F\end{bmatrix}
 \begin{bmatrix}x\\y\end{bmatrix}=0.
 \label{eq:nullpair}
\end{equation}
Because both field bases are linearly independent and $a\ne0$, a nonzero
null vector cannot give $B_Ex=0$: that would also force $y=0$. Thus a single
null-space basis vector suffices. Set $u=B_Ex$, take one exact $t$-th root,
and obtain the other factor by division.

In the special case $t=1$, this is the elementary but powerful criterion
\[
 a\in E_0^{\times}F_0^{\times}
 \quad\Longleftrightarrow\quad E_0\cap aF_0\ne\{0\}.
\]
No factorization of ideals or computation of unit groups is needed for this
\emph{two-fixed-field} membership problem. The unknown factors are
unrestricted field elements, not units or integral elements.

The resulting finite algorithm is to enumerate $d$ increasingly, then the
bounded $t$ in \eqref{eq:tbound}, then all eligible embedded subfields,
and test \eqref{eq:nullpair}. It starts at
\[
 \max\{\lceil\sqrt n\rceil,\ \text{available global lower bounds}\}
\]
and stops no later than $d=n$, where the pair $a\cdot1$ is feasible.

\begin{lstlisting}
g = BuildGaloisData[ap];
r = CompleteTwoFactorDecomposition[ap, g];
r["MaximumDegree"]          (* expected 3 *)
r["OptimalAmongTwoFactors"] (* expected True *)
r["GlobalOptimal"]          (* expected True: lower bound also 3 *)
\end{lstlisting}
The first optimality statement comes from exhaustive testing of the finite
criterion. The second is stronger and is set only when a lower bound valid
for arbitrarily many factors matches the returned degree. These flags must
not be conflated.

\section{A counterexample: optimal factors outside the splitting field}
\label{sec:external}

Consider the positive real number
\begin{equation}
 a=\sqrt{(1+\sqrt2)(1+\sqrt3)}.
 \label{eq:externalalpha}
\end{equation}
It has the evident factorization
\begin{equation}
 a=\underbrace{\sqrt{1+\sqrt2}}_{b}\,
   \underbrace{\sqrt{1+\sqrt3}}_{c}.
 \label{eq:externalfactors}
\end{equation}
The factors have irreducible quartic minimal polynomials
\[
 X^4-2X^2-1,\qquad X^4-2X^2-2.
\]
The minimal polynomial of $a$ is the irreducible degree-eight polynomial
\begin{equation}
 P(X)=X^8-4X^6-16X^4-8X^2+4.
 \label{eq:externalpoly}
\end{equation}
Here $a$ is the \emph{fourth} real root in increasing order, not the first.
Using the positive radical expression avoids any index guess.

We will prove substantially more than failure of the particular factors to
lie in the splitting field: \emph{no product of any number of elements of
that splitting field, all of degree less than eight, can equal $a$}.
Nevertheless \eqref{eq:externalfactors} attains the unrestricted optimum
four.

\subsection{The splitting field has degree sixteen}

Write $r=\sqrt2$, $s=\sqrt3$, $A=1+r$, and $B=1+s$. Set
\[
 B_0=\Q(r,s),\qquad B_1=B_0(i),\qquad L=B_1(a).
\]
The four conjugates of $a^2=AB$ are distinct, so $\Q(a^2)=B_0$.
At some real embeddings of $B_0$, the value of $AB$ is negative; it
therefore cannot be a square in the totally real field $B_0$. It follows
that $[\Q(a):\Q]=8$.

The real elements of $B_1=B_0(i)$ are precisely $B_0$. Since $a$ is real
but not in $B_0$, it is not in $B_1$. Thus $[L:\Q]=16$. All the
conjugates of $a$ lie in $L$ and are
\begin{equation}
 \pm a,\qquad \pm iB/a,\qquad
 \pm irA/a,\qquad \pm r/a.
 \label{eq:allconjugates}
\end{equation}
Conversely, the splitting field contains $\Q(a^2)=B_0$ and the product
$a(iB/a)=iB$, hence also $i$. Therefore $L$ is exactly the splitting
field, rather than merely a convenient larger normal extension.

\subsection{The automorphisms force all small subfields into one place}

Let $z$ fix $B_1$ and send $a$ to $-a$. Define involutions $\sigma$,
$\tau$, and $\kappa$ by
\[
\begin{array}{c|rrrr}
 &r&s&i&a\\\hline
\sigma&-r&s&i&iB/a\\
\tau&r&-s&i&irA/a\\
\kappa&r&s&-i&a
\end{array}
\]
where $\kappa$ is complex conjugation. Direct substitution into the field
relations gives
\[
 z^2=\sigma^2=\tau^2=\kappa^2=1,\qquad z\text{ central},
\]
\[
 [\sigma,\tau]=[\sigma,\kappa]=[\tau,\kappa]=z.
\]
Every group element has a unique form
$\sigma^e\tau^f\kappa^g z^h$, with exponents in $\{0,1\}$, and
\begin{equation}
 (\sigma^e\tau^f\kappa^g z^h)^2=z^{ef+eg+fg}.
 \label{eq:squaregroup}
\end{equation}
The exponent of $\Gal(L/\Q)$ is four.

\begin{lemma}
Every subgroup of $\Gal(L/\Q)$ having order at least four contains $z$.
\end{lemma}
\begin{proof}
Every element of order four squares to $z$. Thus a subgroup of order four
avoiding $z$ would have to be a Klein four group. By
\eqref{eq:squaregroup}, a noncentral involution must lie in one of the
three cosets $\sigma\langle z\rangle$,
$\tau\langle z\rangle$, or $\kappa\langle z\rangle$.
Involutions from different such cosets do not commute; distinct involutions
from the same coset multiply to $z$. Hence no Klein four subgroup avoids
$z$. Every larger subgroup, being a finite 2-group, contains a subgroup of
order four.
\end{proof}

Every subfield of $L$ with degree less than eight has degree one, two, or
four. Its corresponding subgroup has order at least four, so the subfield
is contained in
\[
 L^{\langle z\rangle}=B_1.
\]
Since $a\notin B_1$, no sum or product of elements of these smaller
subfields can equal $a$. This argument allows arbitrarily many such
elements.

\subsection{Exact unrestricted optima}

The quartic factorization gives $\Dtimes(a)\le4$. The group exponent four
and Theorem~\ref{thm:bound} exclude degree caps at most three. Hence
\[
 \boxed{\Dtimes(a)=\Dtwo(a)=4.}
\]
For sums, Lemma~\ref{lem:trace} would move any lower-degree decomposition
into $L$ without increasing the summand degrees. The preceding fixed-field
obstruction rules this out, giving
\[
 \boxed{\Dplus(a)=8.}
\]
Thus the same degree-eight number has an optimal quartic product but no
lower-degree additive decomposition at all.

The two-factor theorem finds \eqref{eq:externalfactors} through
\[
 t=2,\qquad E_0=\Q(\sqrt2),\qquad F_0=\Q(\sqrt3),\qquad u=1+\sqrt2.
\]
Indeed $a^2/u=1+\sqrt3$. Omitting $t>1$ would miss this globally optimal
answer.

\begin{lstlisting}
aext = RootReduce[Sqrt[(1+Sqrt[2]) (1+Sqrt[3])]];
gext = BuildGaloisData[aext];
CompleteTwoFactorDecomposition[aext, gext]
(* expected maximum degree 4; norm exponent 2; globally optimal *)
CompleteSumDecomposition[aext, gext]
(* expected maximum degree 8; globally optimal *)
\end{lstlisting}
The actual factors returned by a null-space basis need not be the positive
radicals displayed above. Any exact branch is permitted by the original
algebraic problem, and the second factor is constructed by division.

\subsection{A specific correction to a published descent claim}

Altamirano's Theorem 3.1 asserts that two lower-degree multiplicative factors
can always be replaced by lower-degree factors inside the Galois closure
\cite[Theorem 3.1, p.~5]{altamirano}. The example above disproves that
assertion. In the printed proof, the proposed replacement uses $a_0/a_1$,
where $a_1$ is the linear coefficient of a relative minimal polynomial;
there is no justification that $a_1\ne0$. Here the relevant polynomial is
$X^2-(1+\sqrt2)$ and the linear coefficient is zero. Retaining the
relative exponent in Theorem~\ref{thm:two} avoids this failure. This
correction concerns that multiplicative descent assertion, not the valid
additive trace reduction.

The quartic product also provides a useful regression test for resultant
inversion:
\begin{lstlisting}
ProductPolynomial[x^4-2 x^2-1, x^4-2 x^2-2, x, z]
(* expected (z^8-4 z^6-16 z^4-8 z^2+4)^2, expanded *)
\end{lstlisting}
Requiring the resultant itself to equal the input minimal polynomial would
reject this genuine factorization.

\section{An extra result for the original product: its best sum uses sextics}
\label{sec:sextics}

The original product $a=uv$ has degree nine and $\Dtimes(a)=3$. Its
additive optimum is different:
\begin{equation}
 \boxed{\Dplus(uv)=6.}
 \label{eq:sexticopt}
\end{equation}
Moreover, restricting summands to $\Q(uv)$ raises the minimum to nine.
This example shows why the additive algorithm needs the whole splitting
field rather than just the input field.

\subsection{A concrete sextic representation}

Let the roots of $p$ be $u_1,u_2,u_3$, and those of $q$ be
$v_1,v_2,v_3$, with $u_1=u$ and $v_1=v$. For $\pi\in S_3$, put
\[
 w_\pi=\sum_{j=1}^{3}u_jv_{\pi(j)}.
\]
There are six such pairing values. The two permutations fixing the real-root
pair give
\[
 w_{\mathrm{id}}+w_{(23)}
 =2u_1v_1+(u_2+u_3)(v_2+v_3)=3u_1v_1.
\]
Both terms on the left are real, because the remaining roots occur in
complex-conjugate pairs. Their common irreducible polynomial is
\begin{align}
 R(X)&=X^6+6X^4-27X^3+9X^2-81X+4\label{eq:sextic}\\
 &=\left(X^3+3X-\frac{27}{2}\right)^2-\frac{713}{4}.\notag
\end{align}
Each of the two cubics obtained by choosing a sign of $\sqrt{713}$ has
strictly positive derivative, so $R$ has exactly two real roots. Consequently
\begin{lstlisting}
w1 = Root[#^6 + 6 #^4 - 27 #^3 + 9 #^2 - 81 # + 4 &, 1];
w2 = Root[#^6 + 6 #^4 - 27 #^3 + 9 #^2 - 81 # + 4 &, 2];
RootReduce[ap - (w1+w2)/3]
(* expected 0 *)
\end{lstlisting}
The summands are $w_1/3$ and $w_2/3$, each still of degree six. If literal
\code{Root} terms with no outside scalar are preferred, use the two real
roots of $R(3X)$.

There is a short algebraic way to check this resolvent. The two real
pairings satisfy
\begin{equation}
 w^2-3uv\,w+(3u^2-3v^2+4)=0.
 \label{eq:pairquadratic}
\end{equation}
Indeed
$(u_2-u_3)^2=-3u^2-4$ and $(v_2-v_3)^2=4-3v^2$.
The polynomial certificates in Section~\ref{sec:examples} express $u,v$
in $\Q(a)$. Substituting them shows that \eqref{eq:pairquadratic} divides
$R(w)$ in $\Q(a)[w]$. This exact reduction is checked independently in
\code{verification/verify.py}.

\subsection{Why no five-degree summands can suffice}

The splitting fields of the two irreducible cubics have groups $S_3$ and
have distinct quadratic subfields $\Q(\sqrt{-31})$ and
$\Q(\sqrt{-23})$. Their intersection, being normal over $\Q$, must be
trivial: the only proper nontrivial normal subfield of an $S_3$ splitting
field is its quadratic subfield. Their compositum therefore has group
\[
 G=S_3\times S_3,\qquad |G|=36.
\]
Let $U=\Span\{u_1,u_2,u_3\}$ and
$V=\Span\{v_1,v_2,v_3\}$. They are the two-dimensional standard
representations, because their only rational root-sum relation is zero.
The products span the four-dimensional irreducible representation
$T=U\otimes V\subset L$, and $a=u_1v_1$ is a nonzero vector of $T$.
Linearly disjoint splitting fields ensure the tensor multiplication map is
injective.

We claim $T^H=0$ for every subgroup $H\le G$ of index at most five.
The possible indices are $1,2,3,4$. Let $A=C_3\times C_3$ be the normal
Sylow 3-subgroup of $G$. For indices one, two, and four, $H$ contains $A$,
which has no invariant vector in $T$. For index three, $|H|=12$,
$H\cap A$ has order three, and the image of $H$ in $G/A=C_2\times C_2$
is the whole quotient. Viewing $A$ as a two-dimensional vector space over
$\mathbb F_3$, its order-three subgroup $H\cap A$ must be invariant under
independent sign changes of the two coordinates. It must therefore be a
coordinate axis. Thus $H$ contains $C_3\times1$ or $1\times C_3$, again
forcing $T^H=0$.

Finite-group representations over $\Q$ admit equivariant projections onto
subrepresentations: average any linear projection over the group. Project
$L$ onto $T$. Every element of every fixed field of degree at most five
projects to zero, whereas $a$ projects to itself. Therefore $a\notin W_5$.
The sextic representation proves $a\in W_6$, establishing
\eqref{eq:sexticopt}.

Finally, $\Q(a)=\Q(u_1,v_1)$ has degree nine, and its stabilizer is
$H_0=S_2\times S_2$, of order four. A subfield of $\Q(a)$ corresponds
to a subgroup containing $H_0$. Its order must be a multiple of four
dividing 36, so the only field degrees are nine, three, and one. All proper
subfields have degree at most three and their elements project to zero in
$T$. Hence even an arbitrarily long lower-degree sum \emph{inside}
$\Q(a)$ is impossible.

The Python verification independently enumerates all 60 subgroups of
$S_3\times S_3$ and reproduces these fixed-space rank calculations. The
proof above does not require trusting that enumeration.

\section{Arbitrarily many product factors: exact search and its limits}
\label{sec:many}

\subsection{Two factors and many factors really differ}

Consider
\[
 \eta=(1+\sqrt2)(1+\sqrt3)(1+\sqrt5).
\]
This number has degree eight, as checked by its exact minimal polynomial in
the verification script. Its displayed three-factor representation gives
$\Dtimes(\eta)=2$. A pair of quadratic factors cannot have degree-eight
product. Nor can a pair with both degrees at most three: if one factor has
degree three, the compositum has degree $3t$ with $1\le t\le3$, which is
not divisible by eight; the remaining cases use only quadratic fields.
Grouping two displayed factors into one degree-four factor gives
\[
 \Dtwo(\eta)=4>2=\Dtimes(\eta).
\]
Thus a two-factor optimum cannot be promoted to a many-factor optimum.
Repeatedly splitting an arbitrary chosen pair also provides no general
proof that the final maximum degree is minimal; it explores only the
factorization tree it happened to choose.

\subsection{A finite dictionary with a precise meaning}

For positive integers $d,h$, let $\mathcal R(d,h)$ be the roots of all
primitive irreducible integer polynomials
\[
 c_0+c_1X+\cdots+c_mX^m,\qquad
 1\le m\le d,\quad |c_j|\le h,\quad 1\le c_m\le h.
\]
Each rational polynomial has a unique primitive integer normalization with
positive leading coefficient, so every algebraic number of degree at most
$d$ occurs once its coefficient-height bound is sufficiently large.
Nonmonic polynomials are necessary: $1/2$, for instance, has primitive
minimal polynomial $2X-1$.

The package enumerates coefficient vectors as base-$(2h+1)$ integers,
rather than allocating all tuples in advance. It tests primitivity and
\code{IrreduciblePolynomialQ}, enumerates every exact root, and removes
canonical duplicates. The raw number of coefficient vectors grows as
\[
 \sum_{m=1}^{d}h(2h+1)^m.
\]
This is intentionally a small-height discovery method, not an efficient
replacement for subfield algorithms at large $d$.

\subsection{Search the prefix and compute the last component exactly}

\code{SearchRootDecomposition[a,op,d,h,r]} searches representations with at
most $r$ components. The first $r-1$ components, when needed, come from
$\mathcal R(d,h)$; the last is an unrestricted exact residual, required
only to have degree at most $d$. For a selected prefix, the residual is
\[
 a-\sum_{j<r}b_j\quad\text{or}\quad
 a\bigg/\prod_{j<r}b_j.
\]
The dictionary therefore does not have to contain every component of a
successful answer. Prefix indices are nondecreasing to avoid permutation
duplicates; repetitions remain allowed. Product searches exclude zero.

At each node the exact residual degree is computed. A node is rejected if
that degree exceeds $d^{\text{remaining slots}}$, or if its largest prime
divisor exceeds $d$. Both tests are necessary conditions, so neither removes
a valid completion. As soon as the residual itself has degree at most $d$,
it becomes the final component. A final \code{RootReduce} equality check
and independent minimal-degree computations validate the candidate.

For the three-quadratic example, the first two factors have coefficient
height two, so the following bounds suffice in the absence of resource
termination:
\begin{lstlisting}
eta = (1+Sqrt[2]) (1+Sqrt[3]) (1+Sqrt[5]);
SearchRootDecomposition[eta, "Product", 2, 2, 3]
(* a successful answer has maximum degree 2 and is globally optimal *)
\end{lstlisting}
The default node limit is 100,000, and the polynomial-enumeration limit is
100,000. These are operational guards, not mathematical height bounds.

\begin{proposition}[What exhaustive bounded search establishes]
With resource limits removed, the search is complete for representations
having at most $r$ components whose first $r-1$ components can be chosen in
$\mathcal R(d,h)$. Letting $h$ and $r$ increase gives a positive-instance
semidecision procedure for the degree cap $d$.
\end{proposition}
\begin{proof}
Every eligible prefix occurs in the finite nondecreasing enumeration after
reordering its commuting components. The residual test recognizes any valid
last component, and the pruning conditions are necessary. For any particular
finite representation at degree cap $d$, its finitely many primitive minimal
polynomials have finite coefficient heights. Some pair of sufficiently
large bounds therefore includes that representation.
\end{proof}

A result \code{NotFoundWithinBounds} means exactly that, not that $d$ is
impossible. Sequentially waiting for an exhaustive search at an infeasible
degree cap can run forever. An anytime implementation can interleave the
finite searches for different caps and height/length budgets, keep the best
verified upper bound, and stop with a global certificate when it matches a
valid lower bound. This article does not supply a proof that such a
procedure terminates with a global certificate for every arbitrary-length
product instance.

\subsection{Why a many-factor norm argument needs more work}

For $a=\prod_i b_i$, take a finite common extension $M/L$ and set
$t_i=[L(b_i):L]$, $T=[M:L]$. Norms give
\begin{equation}
 a^T=\prod_i\Nm_{L(b_i)/L}(b_i)^{T/t_i}.
 \label{eq:manynorms}
\end{equation}
The individual norm values lie in the intersections
$L\cap\Q(b_i)$. Replacing each factor by a $t_i$-th root of its norm
preserves an individual degree upper bound, but the product is only fixed
\emph{up to a root of unity}. Unlike the two-factor construction, one does
not automatically obtain a common relative degree and an exact pairwise
intersection equation.

Related radical replacement, with an explicit root-of-unity factor, is a
central device in Samuels's metric Mahler-measure arguments
\cite[Theorem 2.1]{samuels}. A theorem optimizing Mahler measure is not
therefore a theorem optimizing absolute degree: roots of unity have Mahler
measure one, while their degrees can be significant. For example,
\[
 \zeta_8=\frac{1+i}{\sqrt2}
\]
is a product of degree-two numbers although $\dega(\zeta_8)=4$.
Dropping a root-of-unity correction or inserting it as one supposedly
low-degree component can invalidate a claimed degree bound. Equation
\eqref{eq:manynorms} is useful structure, but not by itself a complete
many-factor minimization algorithm.

\section{Implementation details that affect correctness}
\label{sec:implementation}

\subsection{Constructing exactly the intended Galois field}

The backend begins by computing the rational minimal polynomial of the
input and listing all its roots. It calls
\begin{lstlisting}
common = ToNumberField[roots, All];
\end{lstlisting}
to put them in the \emph{smallest} common field. The \code{All} argument is
important: the documented automatic form need not choose the smallest
common extension \cite{tonumberfield}. A larger normal field would still
support the sum and pair criteria, but its group exponent need not give a
valid lower bound for the original input. The package ties its Galois data
to the exact input and uses the minimal splitting field.

A primitive generator is extracted from one of the resulting
\code{AlgebraicNumber} entries. The generator is then multiplied by the
leading coefficient of its primitive integer minimal polynomial, making it
an algebraic integer without changing the field. This stabilizes
\code{ToNumberField} coordinates: for an integral generator, the documented
representation uses that generator \cite{tonumberfield}. The coordinates
are read with \code{AlgebraicNumberPolynomial}, padded to the field
dimension, and checked to be rational \cite{anpoly}.

Every conjugate of the primitive generator lies in the splitting field.
For each one, powers of its coordinate polynomial modulo the generator's
minimal polynomial produce the corresponding automorphism matrix. The
matrices must be distinct, contain the identity, and be closed under
composition. These checks construct a full multiplication table for the
Galois group without assuming a special built-in subfield or Galois-group
API.

\subsection{Enumerate embedded subgroups, not just conjugacy classes}

Starting with the identity subgroup, the backend repeatedly adjoins an
element and takes closure under the group multiplication table. Finite
positive-power closure already supplies inverses. Duplicate subgroups are
removed, and the process stops when no new subgroup appears.

One must keep all embedded subgroups. Replacing them by representatives up
to conjugacy loses distinct subfields of the chosen copy of $L$ and can
miss a representation of the particular input $a$. Conjugacy compression
is possible only if the missing embeddings are explicitly restored during
the linear-algebra search.

For each subgroup $H$, the computed fixed-space dimension is checked
against $|G|/|H|$. This is a useful independent invariant of the coordinate
construction. A complete data association carries
\code{CompleteSubfieldLattice -> True}; a resource-limited construction
returns a failure instead of such an association.

\subsection{Exact input, explicit exits, and certificate semantics}

The package rejects decimal \code{Real} values. Rationalization of a
numerical approximation is not an exact description of the intended
algebraic number. It asks \code{MinimalPolynomial} for a rational
polynomial and checks the coefficients before using its degree
\cite{minimalpoly}. It does not inspect \code{First[a]} to guess a
polynomial: an exact algebraic input may be a rational, a radical, a
\code{Root}, or a compound expression.

Searches have nested loops and recursion. Wolfram's documented
\code{Return} behavior can leave an inner loop rather than the intended
outer routine \cite{return}. The implementation therefore uses a unique
\code{Catch}/\code{Throw} tag per routine invocation, including each
recursive call. Search-wide resource exits have a separate tag. This keeps
successful returns and failures from being accidentally swallowed by an
inner control construct.

Every returned decomposition is rechecked against the immutable original
input. The \code{Parts} list is the computational result;
\code{Expression} contains an inactive sum or product for presentation.
The flags distinguish exact verification, a two-factor optimum, a global
optimum, and mere failure to find an answer within a budget. Equality and
optimality are separate claims.

\subsection{Public entry points}

\begingroup\small
\begin{longtable}{@{}p{.39\textwidth}p{.57\textwidth}@{}}
\toprule
Function & Meaning and guarantee \\\midrule
\endhead
\code{AlgebraicDegree} & Absolute minimal degree of an exact algebraic input.\\[4pt]
\code{PrimeDegreeBound} & Lower bound valid for every finite sum/product representation.\\[4pt]
\code{PolynomialRoots} & All exact branches of a nonconstant rational polynomial.\\[4pt]
\code{SumPolynomial},\newline\code{ProductPolynomial} & Forward resultant constructions; outputs are annihilating polynomials, not necessarily minimal.\\[4pt]
\code{CheckDecomposition} & Exact equality, component degrees, and a lightweight global certificate when the prime bound is attained.\\[4pt]
\code{PairFromPolynomial} & Candidate-root enumeration plus an exact residual test. Failure is only for that polynomial.\\[4pt]
\code{RootDictionary} & Finite, nonmonic-inclusive coefficient-height enumeration.\\[4pt]
\code{SearchRootDecomposition} & Bounded many-component discovery; no negative global claim from a bounded miss.\\[4pt]
\code{BuildGaloisData} & Explicit splitting field, automorphisms, subgroup lattice, and rational fixed-space bases.\\[4pt]
\code{CompleteSumDecomposition} & Finite criterion for the global additive optimum, subject to successful backend construction.\\[4pt]
\code{CompleteTwoFactorDecomposition} & Finite criterion for the global two-factor optimum, including external factors. Global many-factor optimality requires a matching lower bound.\\
\bottomrule
\end{longtable}
\endgroup

\subsection{Useful optimizations that do not change the mathematics}

The deliberately direct backend can be improved without weakening its
certificates. Cache field-coordinate vectors and multiplication matrices;
store subgroup generators as well as elements; discard linearly redundant
columns during the additive search; and test only degree caps at which the
eligible fixed-field list changes. For pairs, cache the intersections and
use $E_0F_0$ degree bounds to prune field pairs before constructing a large
null-space matrix. Fraction-free or modular linear algebra can reduce
coefficient growth, provided successful reconstructions are checked over
$\Q$.

A more sophisticated number-field backend may compute subfields without
materializing every automorphism matrix. It must still return all relevant
embedded subfields with certified degrees and coordinates. Heuristic
subfield discovery can accelerate positive cases, but cannot certify a
negative degree cap unless completeness is separately established.

\section{Reproducibility and verification status}
\label{sec:verification}

The archive includes the article in source and compiled form, together with
the Wolfram Language package, runnable examples, two native test suites, and
independent verification code with its recorded output. The full package is
also printed in Appendix~\ref{app:code}, so that the article exposes the
implementation rather than hiding it behind function names.

Run the exact independent checks from the archive root with
\begin{lstlisting}[language={}]
python verification/verify.py
\end{lstlisting}
The recorded run used Python 3.13.5 and SymPy 1.14.0. It verifies the two
resultants, irreducibility and real-root counts, rational and polynomial
recovery certificates, the sextic pairing relation, and the degree-eight
external-factor example. It also enumerates the order-36 and order-16
finite groups and checks their subgroup fixed-space obstructions exactly.
All identity and degree acceptance tests use exact arithmetic. Numerical
values are printed only for orientation.

The subgroup counts by index are
\[
\begin{array}{c|rrrrrrrrr}
G=S_3\times S_3,\ [G:H]&1&2&3&4&6&9&12&18&36\\\hline
\#H&1&3&6&1&20&9&4&15&1
\end{array}
\]
and
\[
\begin{array}{c|rrrrr}
\text{external example},\ [G:H]&1&2&4&8&16\\\hline
\#H&1&7&7&7&1.
\end{array}
\]
The finite-group calculations corroborate the proofs rather than replace
them.

\textbf{Native execution limitation.} The available Wolfram evaluator
returned HTTP 404, so it did not execute the package or tests. The
Mathematica outputs in this article are explicitly expected results, not a
claimed kernel transcript. The Python checks validate the algebra and group
calculations, not Wolfram evaluation semantics or end-to-end performance.
The supplied native suites separate basic functionality from expensive
splitting-field tests:
\begin{lstlisting}
TestReport["code/RootDecomposition.wlt"]
TestReport["code/ExtendedTests.wlt"]
\end{lstlisting}
The extended tests include the external quartic factors, the additive
sextic optimum, and the distinction between two and three factors. Running
these is the next implementation-validation step before using the package
as a production simplifier.

To rebuild the article, run \code{latexmk -pdf article.tex}, or run
\code{pdflatex article.tex} twice. The source uses ordinary TeX Live
packages; no font binaries or external figures are required.

\section{Conclusions}

The supplied examples are not merely reducible to cubics: degree three is
the provable global optimum for the indicated operation in each example.
A small polynomial dictionary plus an exact residual calculation discovers
the desired expressions without guessing both cubics simultaneously.

For arbitrary finite sums, trace projection turns the unrestricted problem
in $\overline\Q$ into a finite rational fixed-space problem inside the
splitting field. For two factors, norm descent gives a different finite
criterion: enumerate a bounded radical exponent and test intersections
$E_0\cap a^tF_0$. This retains the auxiliary factors that a naive
splitting-field-only search would discard.

The difference is substantive. The degree-eight number in
Section~\ref{sec:external} has $\Dtimes=4$ but $\Dplus=8$; all
lower-degree elements inside its splitting field lie in a proper subfield.
The original degree-nine product has $\Dtimes=3$ but $\Dplus=6$.
These examples explain why the operation, allowed number of components,
and ambient field must be explicit in both a mathematical theorem and a
Mathematica API.

For products of an arbitrary number of factors, the delivered method is
an exact bounded search with valid global certificates when available, not
a claimed terminating universal minimizer. Keeping that distinction
visible makes the successful answers meaningful: a displayed identity is
verified, and a displayed optimum has an actual lower-bound proof.

\begin{thebibliography}{99}
\bibitem{question}
Vladimir Reshetnikov, \emph{Decomposition of an algebraic number into a sum
or product of algebraic numbers}, Mathematica StackExchange, question
105933, with subsequent answers.
\url{https://mathematica.stackexchange.com/q/105933/7288}.

\bibitem{altamirano}
Christian Altamirano, \emph{A Problem in Algebraic Number Theory}, MIT
UROP+ final paper, 31 August 2017. Mentor: Atticus Christensen; problem
proposed by Bjorn Poonen. See the additive reduction and the multiplicative
Theorem 3.1, discussed critically in Section~\ref{sec:external}.
\url{https://math.mit.edu/research/undergraduate/urop-plus/documents/2017/Altamirano.pdf}.

\bibitem{samuels}
Charles L. Samuels, \emph{The infimum in the metric Mahler measure},
arXiv:1408.4165, 2014. Especially Lemma 2.2 and Theorem 2.1.
\url{https://arxiv.org/abs/1408.4165}.

\bibitem{root}
Wolfram Research, \emph{Root}, Wolfram Language documentation.
\url{https://reference.wolfram.com/language/ref/Root.html}.

\bibitem{rootreduce}
Wolfram Research, \emph{RootReduce}, Wolfram Language documentation.
\url{https://reference.wolfram.com/language/ref/RootReduce.html}.

\bibitem{minimalpoly}
Wolfram Research, \emph{MinimalPolynomial}, Wolfram Language documentation.
\url{https://reference.wolfram.com/language/ref/MinimalPolynomial.html}.

\bibitem{tonumberfield}
Wolfram Research, \emph{ToNumberField}, Wolfram Language documentation.
\url{https://reference.wolfram.com/language/ref/ToNumberField.html}.

\bibitem{anpoly}
Wolfram Research, \emph{AlgebraicNumberPolynomial}, Wolfram Language documentation.
\url{https://reference.wolfram.com/language/ref/AlgebraicNumberPolynomial.html}.

\bibitem{resultant}
Wolfram Research, \emph{Resultant}, Wolfram Language documentation.
\url{https://reference.wolfram.com/language/ref/Resultant.html}.

\bibitem{linearsolve}
Wolfram Research, \emph{LinearSolve}, Wolfram Language documentation.
\url{https://reference.wolfram.com/language/ref/LinearSolve.html}.

\bibitem{return}
Wolfram Research, \emph{Return}, Wolfram Language documentation,
``Possible Issues'' example on nested control constructs.
\url{https://reference.wolfram.com/language/ref/Return.html}.
\end{thebibliography}

\appendix
\clearpage
\section{Complete Wolfram Language reference implementation}
\label{app:code}

This is the exact companion file \code{code/RootDecomposition.wl}. The
mathematical completeness statements and the native-execution limitation
are specified in the main text. The package depends on documented exact
algebraic arithmetic and rational linear algebra, not on external number
field software.

\lstinputlisting[basicstyle=\ttfamily\scriptsize,numbers=left,
 numberstyle=\tiny\color{olive},numbersep=7pt,
 xleftmargin=12pt]{code/RootDecomposition.wl}

\end{document}
